% 395_Messkette_Kanal_De_ch.tex
% Johann Pascher, 24. Okt. 2026
% Die Messkette als Kanal

\providecommand{\BM}{\textbf{[B]}}
\providecommand{\KM}{\textbf{[K]}}
\providecommand{\SM}{\textbf{[S]}}
\providecommand{\XM}{\textbf{[X]}}
\providecommand{\QM}{\textbf{[Q]}}

\chapter*{Die Messkette als Kanal}
\addcontentsline{toc}{chapter}{Die Messkette als Kanal}

\begin{abstract}
Die Nachrichtentechnik bekämpft Echos nicht mehr, sondern misst sie und nutzt sie: Bekannte
Pilotsymbole verraten den Kanal, eine passende Basis macht ihn diagonal, Regeneration verhindert,
dass sich Rauschen aufsummiert, und die Shannon-Grenze sagt, wie viel Information ein verrauschter
Kanal höchstens liefert. Dieses Dokument überträgt die vier Verfahren auf die Frage, ob Rekursionen
Teil der Messwerte sind (Dok.~394). Ergebnisse: Im Verhältnissektor ist der Kanal flach; mit der
Koide-Form als Pilot weichen die Leptonverhältnisse um weniger als $1\,\xi$ ab. Ein modenabhängiger
Kanalfaktor ist damit auf höchstens rund $1\,\%$ des Korrekturschritts $100\,\xi$ beschränkt. Im
SI-Sektor gibt es genau einen Piloten, $\alpha$, und mehr kann es in einer verhältnisbasierten
Theorie nicht geben: Die einzige Beziehung mit einer SI-Einheit ist $\alpha_{\text{SI}}=\xi\,(E_0/1\,\text{MeV})^2$,
$K$ und die MeV-Verankerung sind derselbe Freiheitsgrad. Die Myon-Lebensdauer prüft nicht $K$, sondern
das dimensionslose Verhältnis $v/m_\mu$; die Schranke für den Modenfaktor braucht nur den hergeleiteten
Winkel $2/9$. Für Quantenrechner folgt: Ein gemeinsamer Gatterfaktor wird von der Kalibrierung
aufgenommen und ist prinzipiell nicht messbar; sichtbar wäre nur ein modenabhängiger Anteil von etwa
$10^{-4}$, weit unter der heutigen Dekohärenz. Diskrete Strukturen mit kleinem Nenner überstehen die Messkette, die SI-Korrektur von
$1{,}33\,\%$ lässt nur Nenner bis etwa~$11$ eindeutig. Als Maßstab gegen Zahlenmystik zählt ein
Treffer nur, wenn die Messung mehr Information liefert, als die Formel enthält; die Koide-Parameter
$2/9$ und $\sqrt2$ bestehen das deutlich, ebenso $\alpha^{-1}=3700/27$, wenn man es nach seinen
Bestandteilen $7500\cdot(74/75)/54$ rechnet und nicht als frei gesuchten Bruch.
\end{abstract}

% ============================================================
\section*{Statusmarkierungen}
% ============================================================
\begin{center}
\begin{tabular}{@{}lp{11cm}@{}}
\textbf{[B]} & Algebraisch bewiesen (Umformung ohne Zusatzannahme). \\[2pt]
\textbf{[K]} & Numerisch bestätigt (Prüfskript). \\[2pt]
\textbf{[S]} & Setzung, offen oder Einordnung ohne eigene Rechnung. \\[2pt]
\textbf{[X]} & Durch die Rechnung ausgeschlossen. \\[2pt]
\textbf{[Q]} & Aus der Literatur übernommen, mit Quelle. \\
\end{tabular}
\end{center}
Messwerte: $m_e$, $m_\mu$ und $\alpha$ nach CODATA 2022 \QM, $m_\tau=1776{,}93\pm0{,}09$\,MeV nach PDG
(pdgLive, Ausgabe 2026) \QM.

% ============================================================
\section{Anlass}
% ============================================================

Die Grundform der FFGFT rechnet mit Verhältnissen; eine fraktale Korrektur tritt erst bei der
Übersetzung in SI-Einheiten auf. Dahinter steht die Lesart, dass der Messende selbst innerhalb der
Skalen-Rekursion steht und die Rekursion deshalb Teil jedes Messwerts ist. Dok.~394 hat gezeigt, dass
diese Lesart bisher an genau einer Größe geprüft ist, an $\alpha$, und dass die übrigen Reste andere
Herkunft haben. Offen blieb, wie man die Frage systematisch prüft.

Die Nachrichtentechnik steht seit Jahrzehnten vor derselben Lage: Ein Signal läuft durch einen Kanal,
der es verändert, und der Empfänger steht am Ende dieses Kanals, nicht außerhalb. Wie sie damit
umgeht, liefert vier übertragbare Verfahren. Zwei sind methodisch direkt nutzbar, die Kanalschätzung
mit Pilot und die Shannon-Bilanz; zwei sind Strukturhinweise, die Eigenbasis und die Regeneration.

% ============================================================
\section{Wie die Nachrichtentechnik mit Echos umgeht}
% ============================================================

Ein Funkkanal liefert neben dem direkten Signal verzögerte Kopien, die Echos. Früher galten sie als
Störung; heute werden sie gemessen und als zusätzliche Signalenergie genutzt \QM. Der
\emph{RAKE-Empfänger} (CDMA, UMTS) greift die verzögerten Kopien mit mehreren \enquote{Fingern} einzeln
ab, gleicht ihre Laufzeiten an und addiert sie gewichtet \cite{PriceGreen1958}. \emph{OFDM} (WLAN, LTE,
5G, DVB-T) verteilt das Signal auf viele langsame Unterträger; ein Schutzintervall, das länger ist als
die Echolaufzeiten, macht aus der Faltung mit den Echos eine einfache Phasen- und Amplitudendrehung je
Unterträger \cite{WeinsteinEbert1971}. \emph{Maximum Ratio Combining} und MIMO verwerten mehrere
Ausbreitungswege als parallele Kanäle \cite{Brennan1959}.

Möglich macht das die \emph{Kanalschätzung}: Der Sender schickt Pilotsymbole mit, deren Wert der
Empfänger exakt kennt; aus dem, was ankommt, bestimmt er den Kanal und rechnet ihn heraus. Ein
zweiter Grundsatz ist die \emph{Regeneration}: Jeder Zwischenverstärker entscheidet neu zwischen 0
und~1, deshalb sammelt sich das Rauschen nicht über die Strecke an. Was keines der Verfahren
umgeht, ist die \emph{Shannon-Grenze} $C=B\log_2(1+S/N)$ \cite{Shannon1948}. Echos zu nutzen erhöht das
wirksame Verhältnis $S/N$, und moderne Fehlerkorrektur kommt bis auf Bruchteile eines Dezibels an die
Grenze heran \cite{Berrou1993,Arikan2009}, aber nie darüber. Die Echos verschwinden also nicht; sie
werden gemessen, als Teil des Systems modelliert und verwertet, und die Schranke bleibt stehen.

% ============================================================
\section{Kanalschätzung mit Pilot}
% ============================================================

\subsection{Verfahren und Bedingung}

Übertragen heißt das: Die Messkette ist der Kanal $H$, der den Wert der Grundform $y_{\text{th}}$ in den
Messwert $y_{\text{mess}}=H\,y_{\text{th}}$ überführt. Ein Pilot ist eine Größe, deren Wert in der
Grundform algebraisch exakt feststeht. Aus ihm folgt $H$. Mit diesem $H$ rechnet man die übrigen
Größen zurück. Erklärt ein gemeinsames $H$ mehrere Reste, liegt eine einheitliche Rekursion vor;
braucht jede Größe ihr eigenes $H$, gibt es keine, und man weiß es dann sicher. Die Bedingung dafür
ist streng: Der Pilot muss unabhängig exakt sein. Enthält er selbst einen Modellfehler, schluckt die
Kanalschätzung diesen Fehler, und die Rechnung wird zirkulär.

\subsection{Verhältnissektor: Der Kanal ist flach}

Als Pilot im Verhältnissektor dient die Koide-Form mit $\theta=2/9$ und $d/c=\sqrt2$ (Dok.~353, 367),
\begin{equation}
  \sqrt{m_k}\;\propto\;1+\sqrt2\,\cos\!\Bigl(\tfrac29+\tfrac{2\pi k}{3}\Bigr).
\end{equation}
Sie legt die Verhältnisse $m_\mu/m_e$ und $m_\tau/m_\mu$ ohne Messwert fest. Der Kanalfaktor ist
dann
\begin{equation}
  H_{\mu e}-1=-0{,}074\,\xi,\qquad H_{\tau\mu}-1=-0{,}16\pm0{,}38\,\xi
\end{equation}
\KM, beim Myon-Elektron-Verhältnis bei einer Messunsicherheit von nur $2\times10^{-8}$. Beide Faktoren
liegen unter $1\,\xi$, mehr als hundertmal unter dem SI-Korrekturschritt $100\,\xi$. Im
Verhältnissektor ist der Kanal also flach: Eine Rekursion der Messkette ist dort nicht sichtbar. Das
ist zu erwarten, denn ein gemeinsamer Faktor kürzt sich in Verhältnissen heraus (A040, Dok.~070).

\subsection{Pilot mit Modellfehler: die Leiter}

Nimmt man statt der Koide-Form die Leptonleiter $m_i\propto r_i\,\xi^{p_i}$ als Pilot, liegt sie um
$39{,}10\,\xi$ und $77{,}56\,\xi$ über den gemessenen Verhältnissen \KM; jede Generation verlangt dann
ihren eigenen Kanalfaktor. Das ist aber kein Kanal. Leiter und Koide-Form unterscheiden sich ohne
jeden Messwert um $39{,}02\,\xi$ und $77{,}40\,\xi$ (Dok.~394) \KM: Der scheinbare Kanal ist der
Modellfehler des Piloten. Genau davor warnt die Bedingung aus Abschnitt~3.1. Die fallenden Reste
$\varepsilon_i$ aus Dok.~352 sind damit in der Sprache der Kanalschätzung das Ergebnis eines
ungenauen Piloten.

\subsection{SI-Sektor: ein Pilot, kein Prüfwert}

Im SI-Sektor ist $\alpha$ der einzige Pilot ohne angepassten Faktor. Er schätzt den Kanal sehr genau:
$K=\xi\,m_em_\mu\,\alpha^{-1}=0{,}986501$, also $n=101{,}25$ in $K=1-n\,\xi$, mit einer Unsicherheit
unter $3\times10^{-8}$ \KM; legt man die Energieeinheit über $m_em_\mu=54$ fest, sind es $n=100{,}056$
(Dok.~394) \KM. Ob dasselbe $K$ auch andere SI-Größen trägt, lässt sich nicht prüfen, denn ein zweiter
Pilot fehlt: $E_P$ und $G$ hängen am angepassten Faktor~10 in $v/E_P$ (Dok.~149, 383), $v$ ist eine
Festlegung des Standardmodells (R112). Die Kanalschätzung ist damit möglich, ihre Prüfung nicht. Dass es
einen zweiten SI-Piloten in der FFGFT gar nicht geben kann, zeigt Abschnitt~7.

% ============================================================
\section{Eigenbasis: ein Kanalfaktor je Mode}
% ============================================================

OFDM wechselt in eine Basis, in der der Kanal diagonal wird; aus der Faltung mit Echos wird ein
einfacher Faktor je Unterträger. Die FFGFT hat eine solche Basis bereits: die Moden auf $T^4$ mit ihrer
modengebundenen Zeitskala $\tilde T=1/m$. Eine Messrekursion wäre dort ein eigener Faktor $\kappa_i$ je
Mode, und das ließe generationsabhängige Korrekturen erwarten \SM.

Der Pilot aus Abschnitt~3.2 begrenzt diese Möglichkeit. Ein Unterschied der Kanalfaktoren zweier Moden
erscheint unmittelbar im Verhältnis ihrer Massen. Gilt die Koide-Form exakt, folgt
\begin{equation}
  |\kappa_\mu-\kappa_e|\le0{,}074\,\xi,\qquad |\kappa_\tau-\kappa_\mu|\le0{,}92\,\xi\ (2\sigma)
\end{equation}
\KM. Ein modenabhängiger Anteil der Messrekursion ist damit auf höchstens rund $1\,\%$ des
Korrekturschritts $100\,\xi$ beschränkt. Die SI-Korrektur muss für die drei geladenen Leptonen
praktisch gemeinsam sein; das stützt die einheitliche Regel. Die generationsabhängigen Reste der Leiter
brauchen diese Basis nicht, sie sind der Abstand zur Koide-Form (Abschnitt~3.3). Die Schranke gilt
unter der Annahme, dass die Koide-Form exakt ist; ist sie es nicht, begrenzt die Rechnung nur die Summe
aus Koide-Rest und Modenfaktor \SM. Abschnitt~7 zeigt, dass der Hauptteil der Schranke schon mit dem
hergeleiteten Winkel $2/9$ allein gilt.

Auch die Koide-Form selbst lässt sich frei an die Messwerte anpassen. Mit beiden Parametern frei folgen
$\theta=0{,}2222248\pm0{,}0000063$ und $d/c=1{,}4142089\pm0{,}000011$, also $2/9$ und $\sqrt2$ jeweils
innerhalb einer Standardabweichung, wobei die Unsicherheit allein von der Tau-Masse stammt \KM.

% ============================================================
\section{Regeneration: Diskretes übersteht die Messkette}
% ============================================================

Digitale Übertragung schlägt das Rauschen, weil an jedem Knoten neu entschieden wird. Das deckt sich
mit Dok.~343: Die Diskretheit der FFGFT ist topologisch, sie steckt in Wicklungszahlen und
ganzzahligen Strukturen, nicht in der Dynamik. Übertragen heißt das: Exakte rationale Verhältnisse
und ganze Zahlen lassen sich aus einem verrauschten Messwert wiedergewinnen, solange die Toleranz
kleiner ist als der Abstand zum nächsten Konkurrenten. Die Grenze dafür ist die Farey-Höhe aus
Dok.~343,
\begin{equation}
  Q_c=\frac{\pi}{\sqrt{6\,\varepsilon}} ,
\end{equation}
bis zu der Brüche bei relativer Toleranz $\varepsilon$ eindeutig bleiben:

\begin{center}
\begin{tabular}{@{}lrr@{}}
\toprule
\textbf{Toleranz} & \textbf{$\varepsilon$} & \textbf{$Q_c$}\\
\midrule
SI-Korrektur $100\,\xi$ & $1{,}33\,\%$ & $11{,}1$\\
$m_\tau/m_\mu$ & $5{,}1\times10^{-5}$ & $180$\\
$m_\mu/m_e$ & $2{,}2\times10^{-8}$ & $8700$\\
$\alpha$ & $1{,}5\times10^{-10}$ & $1{,}0\times10^{5}$\\
\bottomrule
\end{tabular}
\end{center}
\KM\ Daraus folgt eine klare Regel für die Darstellung. Wird ein Verhältnis durch die SI-Korrektur
von $1{,}33\,\%$ geführt, bleiben nur Nenner bis etwa~$11$ eindeutig; die Nenner der Grundform $9$,
$5$ und $8$ überstehen das, die $75$ der Rotationszahl nicht. In der Grundform selbst, ohne Korrektur,
sind die Verhältnisse dagegen so genau gemessen, dass auch große Nenner wiedergewinnbar wären.
Fortlaufende Korrekturen sammeln Fehler; exakte Verhältnisse bleiben erhalten. Das stützt die Vorgabe,
die Theorie von der verhältnisbasierten Grundform her aufzuziehen und Korrekturen erst am Ende
anzubringen \SM.

% ============================================================
\section{Shannon-Bilanz gegen Zahlenmystik}
% ============================================================

Eine Messung mit relativer Genauigkeit $\delta$ liefert über eine Konstante, deren Wert vorab nur
auf eine Dekade bekannt ist, rund
\begin{equation}
  I=\log_2\frac{\ln 10}{2\,\delta}\ \text{Bit}.
\end{equation}
Eine Formel enthält selbst Information; Dok.~393 misst sie als Komplexität $C$ in Bit, die mit Zähler,
Nenner, Wurzel und Exponent wächst. Enthält der Raum der möglichen Formeln mehr Bit als die Messung,
trifft zwangsläufig irgendeine Kombination. Daraus folgt ein formales Kriterium: Ein Treffer zählt nur,
wenn der Überschuss $I-C$ deutlich positiv ist. Als $\delta$ gilt die Abweichung vom Messwert, höchstens
aber die Messunsicherheit.

\begin{center}
\begin{tabular}{@{}lrrr@{}}
\toprule
\textbf{Treffer} & \textbf{$I$} & \textbf{$C$} & \textbf{$I-C$}\\
\midrule
$\alpha^{-1}=7400/54$ ($=3700/27$) & $17{,}2$ & $5{,}8$ & $+11{,}5$\\
\quad als frei gesuchter Bruch $3700/27$ & $17{,}2$ & $16{,}6$ & $+0{,}6$\\
$v/E_P=\xi^4/(5\pi)$ & $9{,}0$ & $6{,}3$ & $+2{,}6$\\
$M_W/M_Z=\sqrt{7/9}$ & $10{,}8$ & $5{,}4$ & $+5{,}4$\\
$m_em_\mu$ nahe der ganzen Zahl $54$ & $5{,}8$ & $0$ & $+5{,}8$\\
Koide $\theta=2/9$ & $15{,}3$ & $4{,}2$ & $+11{,}1$\\
Koide $d/c=\sqrt2$ & $17{,}2$ & $2{,}0$ & $+15{,}2$\\
\bottomrule
\end{tabular}
\end{center}
\KM\ Bei $3700/27$ entscheidet, wie man die Formel zählt. Als frei gesuchter Bruch enthält er fast so
viel Information, wie die Abweichung von $7{,}6$\,ppm hergibt, und zählt nicht; so rechnet der
Nulltest aus Dok.~393 ($p\approx0{,}47$). Im Korpus ist $3700/27$ aber nur die Kürzung von
$7500\cdot(74/75)/54$: $7500=1/\xi$ und $74/75$ stammen aus der Grundform, neu ist allein die ganze
Zahl $54$. So gerechnet bleiben über $11$\,Bit Überschuss; unmittelbar geschätzt liegt
$7400/\alpha^{-1}$ mit einer Wahrscheinlichkeit von rund $8\times10^{-4}$ so nahe an einer ganzen Zahl
(Dok.~394) \KM. Voraussetzung ist, dass $74/75$ unabhängig von $\alpha$ feststand (Dok.~133, 300). Die Nähe des
Massenprodukts zu $54$ liefert knapp $6$\,Bit, entsprechend einer Zufallswahrscheinlichkeit von rund
$1{,}7\,\%$ (Dok.~394). Weit vorne stehen die beiden Koide-Parameter mit $11$ und $15$\,Bit Überschuss;
sie sind die stärkste numerische Struktur des Leptonsektors. Dabei ist zu beachten, dass $d/c=\sqrt2$
gleichbedeutend mit der seit 1982 bekannten Koide-Beziehung $Q=2/3$ ist \cite{Koide1982}; der Überschuss
belegt die Struktur der Daten, nicht schon ihre Herleitung aus der FFGFT. Diese steht für $\theta=2/9$
in Dok.~A110 und für die Amplitude in Dok.~353 und 367 (R150).

Das Kriterium ist grob. Es zählt nicht, wie viele Formeln gleicher Komplexität es gibt, und es hängt
vom gewählten Bit-Maß ab \SM. Für einzelne Treffer ersetzt es deshalb nicht den Nulltest aus
Dok.~393, aber es zeigt auf einen Blick, wo sich eine genauere Prüfung lohnt und wo nicht.

% ============================================================
\section{Was sich lösen lässt}
% ============================================================

\subsection{Den zweiten Piloten im SI-Sektor gibt es nicht}

Die einzige Beziehung des Korpus, in der eine SI-Einheit ausdrücklich steht, ist
$\alpha_{\text{SI}}=\xi\,(E_0/1\,\text{MeV})^2$ (Dok.~010, 011, 014, 137, 313, 365, 384, 386). Alle anderen
Aussagen der FFGFT sind reine Zahlen, Verhältnisse zweier Größen gleicher Dimension; in ihnen kürzt sich
ein gemeinsamer Faktor heraus (A040) \BM. Mit $\hbar=c=1$ bleibt eine einzige Dimension, die Energie, und
damit eine einzige Einheit. Der SI-Sektor hat also genau einen Freiheitsgrad: den Faktor, der die Grundform
an das MeV bindet. $K$ und die MeV-Verankerung sind dieselbe Zahl, und die eine Messung $\alpha$ legt sie
fest. Einen zweiten Piloten, der diese Zahl prüfen könnte, kann es in einer verhältnisbasierten Theorie
nicht geben \BM. Die Zuordnung aus Dok.~386, nach der die Korrektur auf der Seite der Einheit sitzt, ist
damit keine bloße Setzung mehr, sondern folgt aus der Struktur; der Wert $74/75$ selbst ist
geometrisch gegeben, als Rotationszahl der eingefrorenen Rekursion mit $1/\xi=7500=100\cdot75$
(Dok.~300, 381, 384) \BM. Offen bleiben nur die Herkunft der Zahl~$100$ (Dok.~133, A040) und der
kleine Rest, den $\alpha$ gegen diesen Wert zeigt (R135) \SM.

Auch die Myon-Lebensdauer ist kein SI-Pilot. Mit $G_F=1/(\sqrt2\,v^2)$ gilt
\begin{equation}
  \frac{\Gamma_\mu}{m_\mu}=\frac{1}{384\,\pi^3}\Bigl(\frac{m_\mu}{v}\Bigr)^4
\end{equation}
bis auf Strahlungskorrekturen \QM; $\tau_\mu$ prüft also nur das dimensionslose Verhältnis $v/m_\mu$. Die
Leiter gibt dafür $5/(16\,\xi)=2343{,}75$, gemessen sind $2330{,}34$: ein Rest von $+43\,\xi$ ohne $K$
und $+145\,\xi$ mit $K$ \KM. Ohne $K$ liegt die Leiter $3{,}4$-mal näher, wie es die Regel verlangt, nach
der $K$ in Verhältnissen keinen Platz hat. Der Rest selbst ist der bekannte Fehler der Leiter: auf $v$
bezogen $+82\,\xi$ beim Elektron, $+43\,\xi$ beim Myon und $-34\,\xi$ beim Tau, mit den Abständen
$39$ und $77\,\xi$ der Leiter zur Koide-Form \KM. $\tau_\mu$ prüft damit die Normierung der Leiter auf
$v$, nicht den Faktor $K$.

\subsection{Die Schranke braucht die Koide-Amplitude kaum}

Die Schranke aus Abschnitt~4 setzte beide Koide-Parameter als exakt voraus. Hergeleitet ist im Korpus
der Winkel $\theta=2/9$ (Dok.~A110); die Begründung der Amplitude $d/c=\sqrt2$ ist offen (R150). Lässt man
die Amplitude frei, schluckt sie einen Teil der Reste, und es bleibt nur eine Kombination der
Modenfaktoren gebunden. Die Amplitude verschiebt aber fast nur das Verhältnis Myon zu Elektron
($\partial\ln(m_\mu/m_e)/\partial a=32{,}6$ gegen $1{,}8$ beim Verhältnis Tau zu Myon) \KM. Die gebundene
Kombination ist deshalb fast genau der Unterschied zwischen Tau und Myon:

\begin{center}
\begin{tabular}{@{}lll@{}}
\toprule
\textbf{exakt angenommen} & \textbf{Tau--Myon} & \textbf{Myon--Elektron}\\
\midrule
$\theta=2/9$ und $d/c=\sqrt2$ & $\le0{,}92\,\xi$ & $\le0{,}074\,\xi$\\
nur $\theta=2/9$ & $\le0{,}91\,\xi$ & nicht gebunden\\
nur $d/c=\sqrt2$ & $\le0{,}92\,\xi$ & nicht gebunden\\
keiner & nicht gebunden & nicht gebunden\\
\bottomrule
\end{tabular}
\end{center}
\KM\ Die Kernaussage, dass ein modenabhängiger Anteil der Messrekursion höchstens rund $1\,\%$ des
Korrekturschritts $100\,\xi$ beträgt, hängt also nur am hergeleiteten Winkel $2/9$. Nur die feinere
Schranke zwischen Myon und Elektron braucht zusätzlich die Amplitude. Sind beide Parameter frei, gibt es
keine Schranke: Zwei Parameter passen zwei Verhältnisse immer exakt an.

% ============================================================
\section{Folgen für Quantenrechner}
% ============================================================

\textbf{Ein gemeinsamer Faktor wird wegkalibriert.} Im Korpus rückt ein Gatter, das um $\pi$
drehen soll, nur um $\pi K$ vor, verkürzt um $100\,\xi=1{,}33\,\%$ (Dok.~034, 175, 230, 231; A070).
R142 hält fest, dass diese Abweichung bei der Pulskalibrierung echter Qubits mit eingestellt wird und
deshalb nicht direkt messbar ist. Die Kanalschätzung gibt dafür den Grund an: Die Kalibrierung ist ein
Pilot. Sie misst, wie weit ein Puls tatsächlich dreht, und stellt ihn so ein, dass das gewünschte
Ergebnis herauskommt. Ein Faktor, der für alle Gatter gleich ist, ist ein einziger Freiheitsgrad, und
den nimmt der Pilot vollständig auf, so wie $\alpha$ den SI-Faktor aufnimmt (Abschnitt~7). Solange $K$
gemeinsam ist, kann ihn kein Quantenrechner sehen, auch nicht mit besserer Hardware \BM.

\textbf{Sichtbar wäre nur ein modenabhängiger Anteil.} Was die Kalibrierung nicht aufnehmen kann, ist ein
Unterschied zwischen Moden, etwa zwischen Qubits verschiedener Frequenz oder zwischen Ein- und
Zwei-Qubit-Gattern. Bei den geladenen Leptonen ist ein solcher Anteil auf höchstens rund $1\,\xi$
beschränkt (Abschnitte~4 und 7), relativ also auf etwa $1{,}2\times10^{-4}$ \KM. Gilt eine ähnliche
Schranke für Qubits, liegt ein mögliches Restsignal in derselben Größenordnung wie die fraktale
Dämpfung der Bell-Größe, $\Delta S/S\approx10^{-5}$ bis $10^{-4}$ (R142). Die Übertragung von Leptonen
auf Qubits ist eine Analogie und nicht hergeleitet \SM. Die heute gemessene Lücke zur Tsirelson-Grenze
$2\sqrt2$ beträgt bei ibm\_kingston $0{,}088$ und gepoolt $0{,}104$ (Dok.~147, 384); sie ist Dekohärenz
und rund $250$- bis $300$-mal größer als ein Restsignal dieser Größe \KM. Damit bleibt der Stand aus
Dok.~384: Operativ liefern FFGFT und Standard-Quantenmechanik dieselben Vorhersagen bis auf eine
$\xi$-Korrektur, die mit heutiger Hardware nicht auflösbar ist.

\textbf{Deterministische Lesart.} Die Ergebnisse stützen sie, ohne sie zu beweisen. Erstens gibt die
offene Bahn aus Dok.~394 einen deterministischen Mechanismus für scheinbaren Zufall: Eine Rotation mit
nicht rationaler Rotationszahl wiederholt ihre Phasen nie exakt und verteilt sie auf lange Sicht
gleichmäßig. Bei einem Vorlauf von $0{,}00056$ Umläufen je Schließungszyklus ist nach rund $450$ Zyklen
eine Viertelumdrehung erreicht, nach rund $1800$ eine volle \KM. Das passt zur Aussage des Korpus,
Rauschen sei ungeordnete Phase (Dok.~384): Jeder Einzelablauf ist festgelegt und wirkt über lange Läufe
dennoch wie Rauschen. Die offene Bahn selbst ist allerdings eine Deutung des Restes von $7{,}6$\,ppm
und nicht hergeleitet \SM. Zweitens überstehen diskrete Ausgänge die Messkette (Abschnitt~5): Ein Bit
ist eine Wicklungszahl (Dok.~343) und lässt sich wie ein digitales Signal regenerieren. Drittens bleibt
die Shannon-Grenze für beide Lesarten dieselbe: Der deterministische Einzelablauf und die
statistische Mittelung über viele Läufe stoßen an denselben Boden (Dok.~343, Satz~G$'$). Die
Bell-Verletzung ist davon nicht berührt; die FFGFT-Messung bleibt kontextuell (Dok.~230).

% ============================================================
\section{Folgerung}
% ============================================================

Die Nachrichtentechnik bekämpft die Rekursion des Kanals nicht, sie misst sie mit Piloten und rechnet
sie heraus; die Shannon-Grenze bleibt dabei bestehen. Übertragen auf die FFGFT ergibt das ein klares
Bild. Im Verhältnissektor ist der Kanal flach: Mit der Koide-Form als Pilot bleibt kein Rest über
$1\,\xi$, und ein modenabhängiger Anteil der Messrekursion ist auf rund $1\,\%$ des Korrekturschritts
begrenzt. Die generationsabhängigen Reste der Leiter sind der Fehler eines ungenauen Piloten, kein
Kanal. Im SI-Sektor gibt es mit $\alpha$ genau einen Piloten, und einen zweiten kann es nicht geben: Die
einzige Beziehung mit einer SI-Einheit legt den einen Freiheitsgrad fest, in dem $K$ und die
MeV-Verankerung zusammenfallen. Die Myon-Lebensdauer prüft die Normierung der Leiter auf $v$, nicht $K$. Diskrete Strukturen mit kleinem Nenner überstehen die Messkette; deshalb gehört die Darstellung in
die Grundform. Und ein Zahlentreffer zählt nur, wenn die Messung mehr Information liefert, als die
Formel enthält. Für Quantenrechner heißt das: Ein gemeinsamer Gatterfaktor verschwindet in der
Kalibrierung, ein modenabhängiger Rest läge bei etwa $10^{-4}$ und damit weit unter der Dekohärenz;
die deterministische Lesart bleibt bis zur Rauschgrenze tragfähig.

Offen bleiben die Herkunft der Zahl~$100$, aus der die Rotationszahl $74/75$ folgt (Dok.~133, 384),
der kleine Rest, den $\alpha$ gegen diesen Wert zeigt (R135), und die Herleitung der Koide-Amplitude (R150); von ihr hängt nur noch die feinere Schranke zwischen Myon und
Elektron ab \SM.

\medskip
\noindent\textbf{Prüfskript:} \path{2/python/Dok395_Skripte/pruef_395_messkette_kanal.py}

\begin{thebibliography}{99}
\bibitem{Shannon1948} C.~E.~Shannon, \textit{A mathematical theory of communication}, Bell System Technical Journal \textbf{27} (1948) 379--423, 623--656.
\bibitem{PriceGreen1958} R.~Price, P.~E.~Green, \textit{A communication technique for multipath channels}, Proceedings of the IRE \textbf{46} (1958) 555--570.
\bibitem{Brennan1959} D.~G.~Brennan, \textit{Linear diversity combining techniques}, Proceedings of the IRE \textbf{47} (1959) 1075--1102.
\bibitem{WeinsteinEbert1971} S.~B.~Weinstein, P.~M.~Ebert, \textit{Data transmission by frequency-division multiplexing using the discrete Fourier transform}, IEEE Transactions on Communication Technology \textbf{19} (1971) 628--634.
\bibitem{Berrou1993} C.~Berrou, A.~Glavieux, P.~Thitimajshima, \textit{Near Shannon limit error-correcting coding and decoding: Turbo-codes}, Proceedings of ICC '93, Geneva (1993) 1064--1070.
\bibitem{Arikan2009} E.~Arıkan, \textit{Channel polarization}, IEEE Transactions on Information Theory \textbf{55} (2009) 3051--3073.
\bibitem{Koide1982} Y.~Koide, Lettere al Nuovo Cimento \textbf{34} (1982) 201; Physics Letters B \textbf{120} (1983) 161.
\end{thebibliography}
